\section{Conclusion}
\label{sec:conclusion}

This paper presented a controlled empirical study of sparse-index gradient accumulation during backpropagation-style AI training workloads. The study focused on standard PyTorch CUDA paths, specifically \texttt{index\_add}, \texttt{scatter\_add}, and \texttt{Embedding.backward}, rather than custom CUDA kernels or full end-to-end model training. The objective was to determine whether these high-level APIs exhibit interchangeable performance behavior under matched sparse-index accumulation workloads, and to relate the observed timing differences to CUDA execution-path evidence.

The results show that high-level API similarity does not imply equivalent CUDA behavior. In the matched benchmark phase, the fastest operation depended on the workload and execution stack. The local GTX1650Ti stack most often favored the embedding backward path, whereas the T4 stacks most often favored \texttt{scatter\_add} under the practical-margin rule. The representative profiling phase further showed that the compared APIs dispatch to distinct CUDA kernel compositions. The \texttt{index\_add} path was dominated by indexed accumulation kernels, the \texttt{scatter\_add} path by scatter/gather-style accumulation kernels, and \texttt{Embedding.backward} by a multi-kernel backward path that included accumulation as well as sorting or segmentation-related components.

These findings support the main conclusion that sparse-index gradient accumulation should be treated as a workload- and stack-dependent systems problem. For embedding-heavy models, recommender-style workloads, graph-like aggregation patterns, and other sparse-feature training settings, selecting an accumulation path only from high-level tensor semantics can be misleading. Focused benchmarking and profiler inspection are useful when indexed gradient accumulation contributes meaningfully to training cost.

The study also defines clear boundaries for this conclusion. The experiments characterize controlled sparse-index accumulation workloads and representative profiling cases; they do not establish a universal fastest PyTorch API, an isolated hardware-causal explanation, or end-to-end training throughput for complete models. The cross-environment comparison is interpreted at the combined hardware/software stack level, since GPU model, operating system, PyTorch version, CUDA version, and runtime conditions differ together.

Future work should extend this characterization to additional GPU generations, PyTorch/CUDA releases, and embedding-heavy training proxies. A natural next step is to evaluate whether the isolated accumulation behavior observed here remains visible inside a complete training step that includes embedding lookup, pooling or feature interaction, a prediction head, loss computation, backward propagation, and optimizer updates. Lower-level profiling with hardware counters could also help explain the microarchitectural causes behind the observed API-dependent behavior.